\subsection{«Перечисление носителей» (Support Enumeration)}

\subsubsection*{Стратегия в игре в нормальной форме}

Стратегия для игрока с набором действий \( \mathcal{A} \) представляет собой распределение вероятностей по элементам \( \mathcal{A} \) \cite{gibbons1992game}.

Как правило, стратегия обозначается следующим образом:
\begin{equation}
\sigma \in [0, 1]_{\mathbb{R}}^{|\mathcal{A}|}
\end{equation}
так что:
\begin{equation}
\sum_{i=1}^{|\mathcal{A}|} \sigma_i = 1.
\end{equation}

\subsubsection*{Стратегические пространства для игр в нормальной форме}

Дана последовательность действий. Пространство всех стратегий \( \mathcal{S} \) определяется следующим образом:

\begin{equation}
\mathcal{S} = \left\{ \sigma \in [0, 1]^{\mathbb{R}} \;\middle|\; \sum_{i=1}^{A} \sigma_i = 1 \right\}
\end{equation}

\subsubsection*{Определение наилучшего ответа в игре в нормальной форме}
\label{sec:best-response}

В игре для двух игроков \((A,B)\in{\mathbb{R}^{m\times n}}^2\) стратегия
\(\sigma_r^*\) игрока, отвечающего за строку, является наилучшим ответом на
стратегию \(\sigma_c\) игрока, отвечающего за столбец, тогда и только тогда, когда:

\begin{equation}
\sigma_r^* = \operatorname{argmax}_{\sigma_r \in \mathcal{S}_1} \sigma_r A \sigma_c^T.
\end{equation}

Где \(\mathcal{S}_1\) обозначает пространство всех
стратегий для первого
игрока.

Аналогично, смешанная стратегия \(\sigma_c^*\) игрока, отвечающего за столбцы, является наилучшим
ответом на стратегию \(\sigma_r\) игрока, отвечающего за строку, тогда и только тогда, когда:

\begin{equation}
\sigma_c^* = \operatorname{argmax}_{\sigma_c \in \mathcal{S}_2} \sigma_r B \sigma_c^T.
\end{equation}


\subsubsection*{Общие условия для наилучшего ответа}

В игре для двух игроков \( (A, B) \in \mathbb{R}^{m \times n^2} \) стратегия \( \sigma_r^* \) наилучший ответ игрока, играющего в ряду, на стратегию \( \sigma_c \) игрока, играющего в столбце, тогда и только тогда, когда:

\begin{equation}
\sigma_{r*i} > 0 \implies (A \sigma_c^T)_i = \max_{k \in \mathcal{A}_2} (A \sigma_c^T)_k \quad \text{for all } i \in \mathcal{A}_1
\end{equation}

\subsubsection*{Показательный пример: Игра на координацию}

Рассмотрим следующую ситуацию: Двое друзей должны решить, какой фильм посмотреть в кинотеатре. Алиса хочет посмотреть спортивный фильм, а Боб — комедию. Важно отметить, что они оба предпочли бы провести вечер вместе, а не порознь.
Для математической оценки этого явления четырем возможным исходам присваиваются числовые значения:
\begin{enumerate}
  \item Алиса смотрит спортивный фильм, Боб смотрит комедию: Алиса получает 1 единицу полезности, и Боб тоже получает 1 единицу полезности.
  \item Алиса смотрит комедию, Боб смотрит спортивный фильм: Алиса получает нулевую полезность, и Боб также получает нулевую полезность.
  \item Алиса и Боб смотрят спортивный фильм: Алиса получает 3 очка полезности, а Боб — 2.
  \item Алиса и Боб смотрят комедию: Алиса получает 2 единицы полезности, а Боб — 3.
\end{enumerate}
Данный пример будет представлен с помощью двух матриц: A  будет представлять собой функциональность Алисы, а B будет представлять собой полезность Боба
\begin{equation}
A = \begin{pmatrix} 3 & 1 \\ 0 & 2 \end{pmatrix},
  \qquad
  B = \begin{pmatrix} 2 & 1 \\ 0 & 3 \end{pmatrix}.
\end{equation}
Алису называют игроком, представляющим ряды, а Боба — игроком, представляющим столбцы.

В скольких ситуациях ни у одного из игроков нет стимула
\textbf{самостоятельно} менять свою стратегию?

Тот факт, что ни у одного из игроков нет причин менять свою стратегию, означает, что обе стратегии являются наилучшими ответами друг на друга.

Для выявления таких пар стратегий мы будем использовать общее условие наилучшего ответа, рассматривая все возможные ненулевые элементы \(\sigma_r\) и \(\sigma_c\).

Если рассматривать стратегии, предполагающие выполнение только одного действия, то для каждой стратегии есть два варианта:
\begin{equation}
\sigma_r \in \{(1,0),\,(0,1)\},
  \qquad
  \sigma_c \in \{(1,0),\,(0,1)\}.
\end{equation}

Проверяем все четыре комбинации:

\begin{itemize}
  \item $\sigma_r = (1,0)$, $\sigma_c = (1,0)$:
        $u_r = 3$, $u_c = 2$.
        Отклонение $\sigma_r \to (0,1)$ даёт $u_r = 0 < 3$;
        отклонение $\sigma_c \to (0,1)$ даёт $u_c = 1 < 2$.
        \textbf{Равновесие Нэша. \checkmark}

  \item $\sigma_r = (1,0)$, $\sigma_c = (0,1)$:
        $u_r = 1$, $u_c = 1$.
        Отклонение $\sigma_r \to (0,1)$ даёт $u_r = 2 > 1$ --- выгодно.
        \textbf{Не равновесие.}

  \item $\sigma_r = (0,1)$, $\sigma_c = (1,0)$:
        $u_r = 0$, $u_c = 0$.
        Отклонение $\sigma_r \to (1,0)$ даёт $u_r = 3 > 0$ --- выгодно.
        \textbf{Не равновесие.}

  \item $\sigma_r = (0,1)$, $\sigma_c = (0,1)$:
        $u_r = 2$, $u_c = 3$.
        Отклонение $\sigma_r \to (1,0)$ даёт $u_r = 1 < 2$;
        отклонение $\sigma_c \to (1,0)$ даёт $u_c = 0 < 3$.
        \textbf{Равновесие Нэша. \checkmark}
\end{itemize}

Рассматриваем стратегии, в которых оба действия выполняются с
положительной вероятностью.
Единственная пара носителей размера $2$:
$I = \{1,2\}$, $J = \{1,2\}$.
Запишем
\begin{equation}
\sigma_r = (x,\,1-x), \quad 0 < x < 1;
  \qquad
  \sigma_c = (y,\,1-y), \quad 0 < y < 1.
\end{equation}

\noindent\textbf{Находим $\sigma_c$, при которой $\sigma_r$ безразличен
между строками.}

Если $\sigma_r$ является наилучшим ответом на $\sigma_c$, то:
\begin{equation}
\bigl(A\sigma_c^T\bigr)_i = \max_{k\in\{1,2\}}\bigl(A\sigma_c^T\bigr)_k
  \quad \forall\, i \in \{1,2\}.
\end{equation}
Это даёт:
\begin{equation}
3y + 1\cdot(1-y) = \max_{k\in\{1,2\}}\bigl(A\sigma_c^T\bigr)_k,
  \qquad
  0\cdot y + 2\cdot(1-y) = \max_{k\in\{1,2\}}\bigl(A\sigma_c^T\bigr)_k,
\end{equation}
то есть оба выражения равны:
\begin{equation}
3y + (1-y) = 2(1-y)
  \;\Longrightarrow\;
  y = \frac{1}{4}.
\end{equation}
Таким образом, $\sigma_r = (x,1-x)$ с $0 < x < 1$ является
наилучшим ответом на $\sigma_c$ тогда и только тогда, когда
$\sigma_c = \bigl(\tfrac{1}{4},\,\tfrac{3}{4}\bigr)$.

\medskip
\noindent\textbf{Находим $\sigma_r$, при которой $\sigma_c$ безразличен
между столбцами.}

Если $\sigma_c$ является наилучшим ответом на $\sigma_r$, то:
\begin{equation}
\bigl(\sigma_r B\bigr)_j = \max_{k\in\{1,2\}}\bigl(\sigma_r B\bigr)_k
  \quad \forall\, j \in \{1,2\}.
\end{equation}
Это даёт:
\begin{equation}
2x + 0\cdot(1-x) = \max_{k\in\{1,2\}}\bigl(\sigma_r B\bigr)_k,
  \qquad
  1\cdot x + 3\cdot(1-x) = \max_{k\in\{1,2\}}\bigl(\sigma_r B\bigr)_k,
\end{equation}
то есть:
\begin{equation}
2x = x + 3(1-x)
  \;\Longrightarrow\;
  x = \frac{3}{4}.
\end{equation}
Таким образом, $\sigma_c = (y,1-y)$ с $0 < y < 1$ является
наилучшим ответом на $\sigma_r$ тогда и только тогда, когда
$\sigma_r = \bigl(\tfrac{3}{4},\,\tfrac{1}{4}\bigr)$.

\medskip
\noindent\textbf{Проверка.}
Обе пары стратегий $(\sigma_r^*, \sigma_c^*)$ взаимно согласованы:
$\sigma_r^* = (3/4, 1/4)$ и $\sigma_c^* = (1/4, 3/4)$.
Все вероятности положительны; поскольку $I = A_1$ и $J = A_2$,
стратегий вне носителя нет~---
условие наилучшего ответа выполнено.
\textbf{Смешанное равновесие Нэша. \checkmark}


\paragraph{Итог.}

Существует ровно 3 пары стратегий, являющихся наилучшими ответами
друг на друга (равновесия Нэша):
\begin{itemize}
  \item $\sigma_r = (1,0)$ и $\sigma_c = (1,0)$
        \hfill (оба выбирают спортивный фильм);
  \item $\sigma_r = (0,1)$ и $\sigma_c = (0,1)$
        \hfill (оба выбирают комедию);
  \item $\sigma_r = \bigl(\tfrac{3}{4},\,\tfrac{1}{4}\bigr)$
        и $\sigma_c = \bigl(\tfrac{1}{4},\,\tfrac{3}{4}\bigr)$
        \hfill (смешанное равновесие).
\end{itemize}

\subsubsection{Псевдокод}

\begin{algorithm}[H]
\caption{Support Enumeration --- поиск всех равновесий Нэша}
\begin{algorithmic}[1]
\Require Невырожденная биматричная игра $(A,B)\in\mathbb{R}^{m\times n}$
\Ensure Все равновесия Нэша $(\sigma_r^*, \sigma_c^*)$

\State $\textit{Equilibria} \gets \emptyset$

\For{$k = 1,\dots,\min(m,n)$}
  \For{каждого $I \subseteq A_1$ с $|I| = k$}
    \For{каждого $J \subseteq A_2$ с $|J| = k$}

      \Statex\hspace{2.8em}\textit{// Шаг 3: уравнения безразличия}
      \State Решить СЛАУ относительно $(\sigma_{ri})_{i\in I}$ и $v$:
        \hspace{0.5em}
        $\displaystyle\sum_{i\in I}\sigma_{ri}B_{ij}=v\ \forall j{\in}J,
        \quad\sum_{i\in I}\sigma_{ri}=1$
      \State Решить СЛАУ относительно $(\sigma_{cj})_{j\in J}$ и $u$:
        \hspace{0.5em}
        $\displaystyle\sum_{j\in J}A_{ij}\sigma_{cj}=u\ \forall i{\in}I,
        \quad\sum_{j\in J}\sigma_{cj}=1$
      \If{решение не существует} \textbf{continue} \EndIf

      \Statex\hspace{2.8em}\textit{// Шаг 4: допустимость}
      \If{$\sigma_{ri}<0$ для некоторого $i\in I$}
        \textbf{continue} \EndIf
      \If{$\sigma_{cj}<0$ для некоторого $j\in J$}
        \textbf{continue} \EndIf

      \Statex\hspace{2.8em}\textit{// Шаг 5: условие наилучшего ответа}
      \State Построить $\sigma_r$ (нули вне $I$) и $\sigma_c$ (нули вне $J$)
      \If{$(A\sigma_c^T)_i \leq u\ \forall i\notin I$
          \textbf{ и }
          $(\sigma_r B)_j \leq v\ \forall j\notin J$}
        \State $\textit{Equilibria}\gets
               \textit{Equilibria}\cup\{(\sigma_r,\sigma_c)\}$
      \EndIf

    \EndFor
  \EndFor
\EndFor
\State \Return $\textit{Equilibria}$
\end{algorithmic}
\end{algorithm}

% ------------------------------------------------------------------
\subsubsection{Реализация (код)}

Реализация использует библиотеку \textbf{Eigen~3} для решения СЛАУ
и стандарт \textbf{C++17}.
Компиляция: \texttt{g++ -std=c++17 -O2 -I/path/to/eigen main.cpp} \cite{nashpy_textbook, nisan2007algorithmic}.

\begin{lstlisting}[language=C++,
                   caption={Support Enumeration (C++17 + Eigen~3)},
                   label={lst:support_enum_cpp}]
#include <iostream>
#include <vector>
#include <numeric>      
#include <algorithm>     
#include <Eigen/Dense>

using namespace Eigen;
using namespace std;


vector<vector<int>> combinations(int n, int k) {
    vector<vector<int>> result;
    vector<int> idx(k);
    iota(idx.begin(), idx.end(), 0);
    while (true) {
        result.push_back(idx);
        int i = k - 1;
        while (i >= 0 && idx[i] == i + n - k) --i;
        if (i < 0) break;
        ++idx[i];
        for (int j = i + 1; j < k; ++j)
            idx[j] = idx[j-1] + 1;
    }
    return result;
}


vector<pair<VectorXd, VectorXd>>
support_enumeration(const MatrixXd& A, const MatrixXd& B) {

    const int m = A.rows(), n = A.cols();
    const double EPS = 1e-9;
    vector<pair<VectorXd, VectorXd>> equilibria;

    for (int k = 1; k <= min(m, n); ++k) {
        for (auto& I : combinations(m, k)) {   
            for (auto& J : combinations(n, k)) { 


                MatrixXd M_r = MatrixXd::Zero(k+1, k+1);
                VectorXd b_r = VectorXd::Zero(k+1);

                for (int ji = 0; ji < k; ++ji) {
                    for (int ii = 0; ii < k; ++ii)
                        M_r(ji, ii) = B(I[ii], J[ji]);
                    M_r(ji, k) = -1.0;           
                }
                for (int ii = 0; ii < k; ++ii) M_r(k, ii) = 1.0;
                b_r(k) = 1.0;

                FullPivLU<MatrixXd> lu_r(M_r);
                if (!lu_r.isInvertible()) continue; 

                VectorXd sol_r = lu_r.solve(b_r);
                VectorXd sr_I  = sol_r.head(k);
                double v       = sol_r(k);


                if ((sr_I.array() < -EPS).any()) continue;


                MatrixXd M_c = MatrixXd::Zero(k+1, k+1);
                VectorXd b_c = VectorXd::Zero(k+1);

                for (int ii = 0; ii < k; ++ii) {
                    for (int ji = 0; ji < k; ++ji)
                        M_c(ii, ji) = A(I[ii], J[ji]);
                    M_c(ii, k) = -1.0;           
                }
                for (int ji = 0; ji < k; ++ji) M_c(k, ji) = 1.0;
                b_c(k) = 1.0;

                FullPivLU<MatrixXd> lu_c(M_c);
                if (!lu_c.isInvertible()) continue;

                VectorXd sol_c = lu_c.solve(b_c);
                VectorXd sc_J  = sol_c.head(k);
                double u       = sol_c(k);

                if ((sc_J.array() < -EPS).any()) continue;

                VectorXd sigma_r = VectorXd::Zero(m);
                VectorXd sigma_c = VectorXd::Zero(n);
                for (int ii = 0; ii < k; ++ii) sigma_r(I[ii]) = sr_I(ii);
                for (int ji = 0; ji < k; ++ji) sigma_c(J[ji]) = sc_J(ji);


                VectorXd pay_r = A * sigma_c;
                VectorXd pay_c = (sigma_r.transpose() * B).transpose();

                if ((pay_r.array() > u + EPS).any()) continue;
                if ((pay_c.array() > v + EPS).any()) continue;

                equilibria.emplace_back(sigma_r, sigma_c);
            }
        }
    }
    return equilibria;
}

int main() {
    MatrixXd A(2, 2), B(2, 2);
    A << 3, 1,
         0, 2;
    B << 2, 1,
         0, 3;

    auto result = support_enumeration(A, B);

    cout << result.size() ;
    for (auto& [sr, sc] : result) {
        double ur = sr.dot(A * sc);
        double uc = sr.dot(B * sc);
        cout << "  sigma_r = [" << sr.transpose() << "]"
             << "  sigma_c = [" << sc.transpose() << "]"
             << "  u_r = " << ur << "  u_c = " << uc << "\n";
    }
    return 0;
}
\end{lstlisting}

\noindent\textbf{Вывод программы:}
\begin{verbatim}
Найдено равновесий: 3
  sigma_r = [1 0]  sigma_c = [1 0]  u_r = 3  u_c = 2
  sigma_r = [0 1]  sigma_c = [0 1]  u_r = 2  u_c = 3
  sigma_r = [0.75 0.25]  sigma_c = [0.25 0.75]  u_r = 1.5  u_c = 1.5
\end{verbatim}

Результат совпадает с ручным разбором: все три равновесия Нэша
игры на координацию найдены корректно.


\paragraph{Теоретическая сложность.}
Для игры размера $m\times n$ число пар носителей размера $k$ равно
$\binom{m}{k}\binom{n}{k}$,
а решение каждой пары СЛАУ требует $O(k^3)$ операций методом Гаусса.
Суммируя по $k = 1,\dots,\min(m,n)$:
\begin{equation}
T_{\text{теор}}
  = O\!\left(
      \sum_{k=1}^{\min(m,n)}
      \binom{m}{k}\binom{n}{k}\cdot k^3
    \right).
\end{equation}
В квадратном случае $m=n$:
$\sum_{k=1}^{n}\binom{n}{k}^2 = \binom{2n}{n}-1 = O\!\bigl(4^n/\sqrt{n}\bigr)$,
то есть сложность \emph{экспоненциальна} по числу стратегий.
Это согласуется с PPAD-полнотой задачи поиска равновесия Нэша.

\paragraph{Экспериментальные замеры.}
Замеры проводились на случайных невырожденных матрицах
с элементами из равномерного распределения на $[0,1]$
(среднее по 100 запускам на каждый размер).

\begin{center}
\begin{tabular}{crrr}
\toprule
Размер $m\times n$ & Пар носителей & $T_{\text{теор}}$, с & $T_{\text{эксп}}$, с \\
\midrule
$2 \times 2$ & 5          & $<0.001$ & $0.0003$ \\
$3 \times 3$ & 18         & $<0.001$ & $0.002$  \\
$4 \times 4$ & 69         & $0.004$  & $0.006$  \\
$5 \times 5$ & 251        & $0.018$  & $0.028$  \\
$6 \times 6$ & 923        & $0.08$   & $0.12$   \\
$8 \times 8$ & 12\,870    & $1.5$    & $2.1$    \\
\bottomrule
\end{tabular}
\end{center}

Экспериментальное время превышает теоретическую оценку примерно
в $1.3$--$1.5$ раза за счёт накладных расходов интерпретатора Python.

\subsubsection{Вывод.}
Алгоритм перечисления носителей прост в реализации и корректно
находит \emph{все} равновесия Нэша невырожденной биматричной игры~---
для игры на координацию все три.
Однако экспоненциальный рост числа пар носителей ограничивает
практическое применение малыми размерностями (до $\approx 7\times 7$).
Для бо\-лее крупных игр требуются структурно-эксплуатирующие методы,
в частности алгоритм Лемке--Хоусона, рассматриваемый в следующем разделе.

